% ============================================================
% CHAPTER 15: FUTURE WORK
% MCDA 5585 / 5586 Major Project Report
% Bhavik Kantilal Bhagat | A00494758 | Saint Mary's University
% ============================================================

\chapter{Future Work}
\label{chap:future_work}

The work I delivered in this co-op placement establishes a monitoring foundation ---
not an endpoint. I have already defined several planned extensions in the JIRA
backlog, scoped for implementation during the Mitacs BSI renewal
(Sep~1 -- Dec~31, 2026).

\paragraph{Distributed Tracing with X-Ray.}
The most technically compelling near-term extension is the integration of
AWS~X-Ray distributed tracing into the IA Lambda fleet using the
\texttt{aws-lambda-powertools[tracer]} library. X-Ray would provide
end-to-end request traces across Lambda invocations, SQS message dispatch,
and downstream HTTP calls --- the ``service map'' view that CloudWatch metrics
alone cannot produce. At the 5\% sampling rate recommended for
production workloads, the additional cost is estimated at under \$2 per month
for the current fleet size, making this an unusually cost-effective
observability enhancement. A Grafana X-Ray data source plugin is available
for AMG, meaning trace data would flow directly into the existing workspace
without requiring a separate visualization tool.

\subsubsection{Auto-Healing Agent}
\label{sec:autohealing_agent}

The alarm infrastructure I delivered closes the detection gap --- the 2--3 hour
window between failure onset and SRE awareness that the Meridian Incident
exposed. But detection without automated remediation still leaves a response
gap: a human must receive the SNS notification, triage the alarm, identify
the correct recovery action, and execute it. For the most common and
well-understood failure modes on the IA platform, this human step is
unnecessary. The \textbf{auto-healing agent} eliminates it.

\paragraph{Concept.}
The auto-healing agent is an event-driven remediation layer that sits between
the CloudWatch alarm infrastructure and the AWS services being monitored.
Each alarm SNS message triggers a Lambda ``healer'' function. The healer
evaluates the alarm payload, classifies the failure type, selects the
appropriate recovery action from a configurable remediation catalogue, and
executes it --- logging the entire sequence via the IA Event Store with
structured PROCESS\_EVENT codes so that every automated intervention is
fully auditable.

\paragraph{Architecture.}
The proposed architecture connects directly to the existing dashboard alarm
infrastructure without requiring changes to the Grafana dashboard or
CloudFormation alarm stacks:

\begin{enumerate}
    \item A CloudWatch Alarm breaches its threshold (e.g., ECS task desired
          count $\neq$ running count on a Meridian service).
    \item The alarm fires an SNS message to the \texttt{ia-sre-autohealing-topic}
          SNS topic (a new topic, parallel to the existing notification topic).
    \item The SNS topic triggers the \textbf{Healer Lambda}, which receives
          the full alarm payload including the \texttt{AlarmName},
          \texttt{NewStateValue}, \texttt{MetricName}, and any available
          dimensions (\texttt{ClusterName}, \texttt{ServiceName},
          \texttt{FunctionName}).
    \item The Healer Lambda looks up the alarm name in a
          \textbf{remediation catalogue} --- a DynamoDB table mapping alarm
          patterns to recovery action functions.
    \item The matched recovery action is invoked: an AWS API call, a Step
          Functions state machine, or a queued message to a downstream
          remediation worker.
    \item The Healer Lambda logs the intervention to the IA Event Store
          with a \texttt{PROJ-BE-AUTOHEALING} code, including the alarm
          name, recovery action taken, and outcome status.
    \item If the recovery action fails or the alarm re-fires within a
          configurable cool-down window, the Healer escalates by publishing
          to the existing notification SNS topic to alert the SRE team
          directly.
\end{enumerate}

\paragraph{Remediation Catalogue (Initial Scope).}
The first phase of the auto-healing agent targets the failure modes observed or anticipated
from the four support rotation weeks:

\begin{table}[ht]
\centering
\caption{Auto-Healing Agent --- Initial Remediation Catalogue}
\label{tab:autohealing_catalogue}
\rowcolors{2}{tablerowA}{tablerowB}
\begin{tabularx}{\textwidth}{p{0.22\textwidth}p{0.48\textwidth}p{0.22\textwidth}}
\hline
\rowcolor{smubrand}\textcolor{white}{\textbf{Alarm / Trigger}} & \textcolor{white}{\textbf{Recovery Action}} & \textcolor{white}{\textbf{Escalation if Fails}} \\
\hline
ECS task desired $\neq$ running & \texttt{ecs:UpdateService} (force new deployment) & SNS alert to SRE \\
MSK consumer lag $>$ 10{,}000 messages & Scale ECS consumer service (+1 task) & SNS + JIRA ticket (auto-created via JIRA API) \\
Lambda error rate $>$ 1\% over 5 min & Trigger Step Functions diagnostic workflow; capture last 50 error log lines to S3 & SNS alert with S3 log link \\
SQS queue age $>$ 1 hour & Invoke queue-processor Lambda with \texttt{force=true} parameter & SNS alert \\
ECS CPU $>$ 80\% sustained 5 min & Trigger ECS service auto-scaling policy (\texttt{RegisterScalableTarget}) & SNS if already at max task count \\
RPA bot missing heartbeat (future) & Call UiPath Orchestrator API to restart robot job & SNS + support ticket \\
\hline
\end{tabularx}
\end{table}

\paragraph{Design Principles.}

Three design principles constrain the auto-healing agent to avoid introducing
new failure modes:

\begin{itemize}
    \item \textbf{Idempotency.} Every recovery action must be safe to invoke
          multiple times. The DynamoDB remediation catalogue tracks the last
          invocation time per alarm; the Healer Lambda checks this before
          acting and skips duplicate invocations within the cool-down window.

    \item \textbf{Bounded autonomy.} The Healer acts on a whitelist of
          known-safe recovery actions. Any alarm with no matching catalogue
          entry falls through to standard SNS notification --- the agent
          does not improvise. New recovery actions are added only after
          review and approval, following the same PR process as
          CloudFormation template changes.

    \item \textbf{Full auditability.} Every automated intervention is
          logged to the IA Event Store with a structured event code
          (\texttt{PROJ-I-AUTOHEALING}), the alarm trigger name, the
          action executed, the outcome status, and a link to the AWS
          resource modified. The Event Store trail ensures that automated
          interventions are visible in the same audit log as manual ones.
\end{itemize}

\paragraph{Relation to the PoC.}
The auto-healing agent is the direct successor to the alarm-and-notify
posture delivered by this co-op placement. The detection layer (CloudWatch
alarms + SNS) was built first because a remediation layer without a
reliable detection layer is fragile --- automated actions triggered by
false alarms cause more harm than the failures they prevent. Now that the
detection infrastructure has been in production for several months, the
signal quality of the alarms is validated, and the false-positive rate
is well-characterized. The auto-healing agent can therefore be implemented against a
tested alarm baseline rather than a hypothetical one.

This progression --- detect first, then automate remediation --- is the
correct sequence for building reliable auto-healing in a regulated financial
services environment, where an incorrectly triggered service restart during
a VPL pricing window or a CAIS Deadline processing run would itself
constitute a production incident.

\paragraph{RPA Bot CloudWatch Integration.}
Two follow-on stories in my JIRA backlog propose bridging the Blue Prism
(BPM) and UiPath Cloud operational metrics into CloudWatch via a polling
Lambda function. The Lambda would query the Blue Prism API and UiPath
Orchestrator API on a configurable schedule, emit custom CloudWatch metrics
for queue depth, bot session count, and process error rates, and feed those
metrics into the existing Grafana dashboard infrastructure. This would close
the last remaining visibility gap: the current dashboard covers the AWS
infrastructure layer comprehensively but has no visibility into the RPA
execution layer above it.

\paragraph{Error Budget Burn-Rate Alerting.}
The current alarm model fires on threshold breaches without distinguishing
between a momentary spike and a sustained reliability regression. Multi-window
burn-rate alerting --- as described in the Google SRE Workbook --- uses
1-hour and 6-hour Metric Math expressions to detect when the error budget is
being consumed faster than the replenishment rate. Implementing this pattern
in CloudWatch Metric Math would elevate the existing alarm infrastructure from
threshold-based detection to SLO-aware alerting, enabling the team to
distinguish incidents that require immediate response from those that are
within the acceptable error budget envelope.

\paragraph{Tag Compliance Automation.}
The AppName tagging obligation I identified in Section~\ref{sec:challenges-infra} currently
relies on a reactive CloudWatch alarm: the alarm fires after a non-compliant
resource has already been deployed. A proactive alternative would query the
inventory API on a schedule, compare the returned resource list against an
expected tag manifest, and open a JIRA ticket automatically for any resource
found to be non-compliant. This would shift tag compliance from a detection
activity to a governance workflow.

\paragraph{Mitacs BSI Renewal Scope.}
The approved Mitacs BSI renewal (Sep~1 -- Dec~31, 2026) will extend
my industrial collaboration with a focus on three themes: X-Ray tracing
integration, IaC hardening and drift detection for the CloudFormation
stacks I delivered in this period, and capacity forecasting using CloudWatch
anomaly detection models trained on the historical metric data now being
accumulated by the deployed dashboards.


% Keep all floats within this chapter
\FloatBarrier
